\begin{pastexam}[label={px:gw3-114-ds-04}]{關務三等資訊處理 114 年 第四題}
（申論；25 分）

請逐步寫出下列使用遞迴函式的呼叫與輸出過程。（25 分）
\begin{minted}[fontsize=\small, style=perldoc]{cpp}
#include <iostream>
using namespace std;
int A(int n, int c = 1) {
     if (n == 0)
           return c + 1;
     return A(n - 2, c * n);
}
int main() {
     cout << A(6) << endl;
     return 0;
}
\end{minted}
\end{pastexam}

\begin{pastsoln}
\texttt{A(6)} 只給了一個引數，第二個參數 \texttt{c} 用預設值 $1$，所以實際呼叫的是 \texttt{A(6, 1)}。 每次
$n \ne 0$ 就呼叫 \texttt{A(n - 2, c * n)}：$n$ 每次減 $2$，$c$ 累乘目前的 $n$；$n = 0$ 時傳回 $c + 1$。

\textbf{呼叫過程}（往下呼叫）：
\begin{center}
\begin{tabular}{clll}
\toprule
層 & 呼叫 & $n = 0$？ & 動作 \\
\midrule
1 & \texttt{A(6, 1)} & 否 & 呼叫 \texttt{A(6-2, 1*6)} $=$ \texttt{A(4, 6)} \\
2 & \texttt{A(4, 6)} & 否 & 呼叫 \texttt{A(4-2, 6*4)} $=$ \texttt{A(2, 24)} \\
3 & \texttt{A(2, 24)} & 否 & 呼叫 \texttt{A(2-2, 24*2)} $=$ \texttt{A(0, 48)} \\
4 & \texttt{A(0, 48)} & 是 & 傳回 $48 + 1 = 49$ \\
\bottomrule
\end{tabular}
\end{center}

\textbf{返回過程}（往上傳回）：第 4 層傳回 $49$ 給第 3 層；第 3 層的 \texttt{return A(0, 48);} 直接把
$49$ 傳回給第 2 層；同樣第 2 層傳回 $49$ 給第 1 層，第 1 層傳回 $49$ 給 \texttt{main}。 呼叫堆疊的變化：
\begin{gather*}
\texttt{main} \to \texttt{A(6,1)} \to \texttt{A(4,6)} \to \texttt{A(2,24)} \to \texttt{A(0,48)} \\
\texttt{A(0,48)} \xrightarrow{\ 49\ } \texttt{A(2,24)} \xrightarrow{\ 49\ } \texttt{A(4,6)}
\xrightarrow{\ 49\ } \texttt{A(6,1)} \xrightarrow{\ 49\ } \texttt{main}
\end{gather*}

\textbf{輸出}：\texttt{main} 印出 \texttt{49} 並換行。
\begin{minted}[fontsize=\small]{text}
49
\end{minted}

\textbf{說明}：$c$ 一路累積 $6 \times 4 \times 2 = 48$，所以 \texttt{A(n)}（$n$ 為偶數）算的是
$n!! + 1$（$n!! = n(n-2)(n-4)\cdots 2$），$6!! + 1 = 48 + 1 = 49$。 這個函式的遞迴呼叫是最後一個動作
（尾遞迴），回來後不必再做任何計算，可以直接改寫成迴圈（資結書第~6.4~節）。 另外注意：若 $n$ 是奇數，
$n$ 會跳過 $0$ 變成負數，永遠不會停，最後堆疊溢位；題目的 $n = 6$ 沒有這個問題。
\end{pastsoln}
